% Prepared for Edward Baker, 4 November 2026.
% Consolidated with substantial LLM assistance: gpt-6-astra, effort ultra.
% Manuscript fragment; requires amsmath, amssymb, amsthm and the proposition environment.
% Local evidence paths are collected by the main manuscript in sec:provenance.
\section{Green transfer, complete boundary terms, and finite comparisons}
\label{sec:green-comparison}

The following calculation retains the undifferentiated theta primitive and
its inner boundary.  It produces a quantitative finite comparison, but
positivity of the primitive alone does not sign the resulting determinants.
The comparison and its limitations consolidate RGD Notes~01, 05 and~07;
their exact paths and evidence relationships are recorded in
Appendix~\ref{sec:provenance}.  Throughout this section the physical heat
time is zero, and all negative coefficient or moment indices mean zero.

\subsection{Normalization and the complete boundary expansion}
Put
\begin{align}
 p_n(u)&=\exp(u-\pi n^2e^{4u}), & p(u)&=\sum_{n\ge1}p_n(u),\\
 \Phi(u)&=\sum_{n\ge1}\bigl(2\pi^2n^4e^{9u}-3\pi n^2e^{5u}\bigr)
                    e^{-\pi n^2e^{4u}}, &
 m_j(u)&=\frac{u^{2j}}{(2j)!}\quad(j\ge0),\\
 c_j&=\int_0^\infty\Phi(u)m_j(u)\,du, &
 d_j&=\int_0^\infty p(u)m_j(u)\,du.
\end{align}
We use ordinary coefficients, so
\begin{equation}
 C(z)=\sum_{j\ge0}c_jz^j
 =\int_0^\infty\Phi(u)\cosh(u\sqrt z)\,du
 =\frac18\xi\!\left(\frac{1+\sqrt z}{2}\right),
 \qquad C(1)=b:=\frac1{16}.
 \label{eq:green-C}
\end{equation}
The even power series makes the square-root notation single-valued.
Define
\begin{equation}
 D_{r,k}=\det[c_{k+j-i}]_{0\le i,j<r},\qquad
 D_{0,k}=1,\qquad r\ge1,\ k\ge0.
 \label{eq:green-D}
\end{equation}

For clarity, the radial origin of the transfer can be checked directly.
With \(R=e^{2u}\) and
\(A(R)=\sum_{n\ge1}e^{-\pi n^2R^2}\), one has
\(p(u)=R^{1/2}A(R)\) and
\begin{equation}
 \Phi(u)=\frac12R^{5/2}\left(\partial_R^2+\frac2R\partial_R\right)A(R),
 \qquad p''-p=8\Phi.
\end{equation}
Theta reflection gives
\(1+2A(R)=R^{-1}[1+2A(1/R)]\).  Differentiating at \(R=1\)
yields \(A'(1)=-A(1)/2-1/4\), hence \(p'(0)=-1/2\).
The complete source tail makes every fixed polynomial derivative transfer
below absolutely convergent and removes the outer boundary.

\begin{proposition}[Green transfer with the inner boundary]
For every polynomial \(f\),
\begin{equation}
 \int_0^\infty\Phi f
 =b f(0)+\frac{p(0)}8f'(0)
   +\frac18\int_0^\infty p(f''-f).
 \label{eq:green-transfer}
\end{equation}
In particular,
\begin{equation}
 c_j=b\delta_{j0}+\frac18(d_{j-1}-d_j),\qquad
 c_0=b-d_0/8.
 \label{eq:green-coefficients}
\end{equation}
\end{proposition}
\begin{proof}
Integrate \((p''-p)f/8\) twice by parts.  The inner contribution is
\([-p'(0)f(0)+p(0)f'(0)]/8\).  Now use
\(m_j''=m_{j-1}\), \(m_j'(0)=0\), and \(m_j(0)=\delta_{j0}\).
\end{proof}
The tangent inequality used below also gives
\begin{equation}
 0<d_0\le\frac{e^{-\pi}}{(4\pi-1)(1-e^{-3\pi})}<\frac14,
 \qquad \frac1{32}<c_0<b.
\end{equation}
Indeed \(n^2\ge1+3(n-1)\) controls the complete sum; the last coarse
bound follows already from \(\pi>3\), \(e^3>8\) and \(e^9>2\).

The determinant version includes every lower-dimensional boundary term.
Let
\begin{equation}
 \mathcal T_{r,k}(\boldsymbol u)
 =\frac1{r!}\sum_{\sigma\in S_r}
   \det[m_{k+j-i}(u_{\sigma(i)})]_{0\le i,j<r},
\end{equation}
and, for \(0\le q\le r\), put
\begin{equation}
 W_{r,k}^{(q)}(u_0,\ldots,u_{q-1})
 =\left[\left(\prod_{i=0}^{q-1}(\partial_{u_i}^2-1)\right)
       \mathcal T_{r,k}(\boldsymbol u)\right]_{u_q=\cdots=u_{r-1}=0}.
\end{equation}
An empty integral denotes evaluation of its scalar integrand.
\begin{proposition}[Complete Green determinant identity]
For every \(r\ge1\) and \(k\ge0\),
\begin{equation}
 D_{r,k}=\sum_{q=0}^r\binom rq b^{r-q}8^{-q}
 \int_{[0,\infty)^q}\prod_{i=0}^{q-1}p(u_i)\,
                    W_{r,k}^{(q)}(\boldsymbol u)\,d\boldsymbol u.
 \label{eq:green-boundary-determinant}
\end{equation}
Moreover, \(W_{r,k}^{(q)}\equiv0\) if \(q<\min(k,r)\).
Thus only the full-volume term remains when \(k\ge r\), while
\(D_{r,0}=c_0^r\).
\end{proposition}
\begin{proof}
Expanding the row determinant and integrating against
\(\prod_i\Phi(u_i)\) gives \(D_{r,k}\); symmetrizing does not change the
integral.  Apply \eqref{eq:green-transfer} in each coordinate.
Coordinatewise evenness removes every first-derivative trace, and symmetry
makes the \(\binom rq\) terms with \(q\) volume coordinates equal.
An endpoint row has entries \(\delta_{k+j-i,0}\).  At most
\(\max(r-k,0)\) such rows can be nonzero, proving the asserted vanishing.
The triangular coefficient matrix at \(k=0\) proves the final identity.
\end{proof}
There is no further factorial in \eqref{eq:green-boundary-determinant}:
\(\mathcal T\) was defined by an average.  The probes and primitive
moment differences can have either sign.  In particular, neither
\(p>0\) nor the pure-volume representation is a positivity theorem.

\subsection{Preserving the boundary before approximating the source}
Let
\begin{equation}
 x_j=\frac18\int_0^\infty p_1m_j,\qquad
 y_j=\frac18\int_0^\infty\sum_{n\ge2}p_nm_j,\qquad
 \widehat c_j=b\delta_{j0}+x_{j-1}-x_j.
\end{equation}
Set \(a=16\pi-1\).  The tangent bound \(e^{4u}\ge1+4u\) and
\(n^2\ge4+5(n-2)\) for \(n\ge2\) give
\begin{equation}
 0<y_j\le E_j:=
 \frac{e^{-4\pi}}{8(1-e^{-5\pi})a^{2j+1}}.
 \label{eq:green-mode-tail}
\end{equation}
Indeed, the exponential tangent bounds the \(n\)-th summand by
\(e^{-\pi n^2}e^{-(4\pi n^2-1)u}\), and the normalized moment
integral equals \((4\pi n^2-1)^{-2j-1}\); summing its geometric
majorant proves \eqref{eq:green-mode-tail}.

This reference has a precise positive-measure description, but no
all-order Toeplitz positivity is assumed for it.  Write
\(\beta_n=(4\pi n^2-1)e^{-\pi n^2}/8\) and
\(\tau=b-\beta_1=\sum_{n\ge2}\beta_n>0\).
Individual-mode Green transfer shows that \((\widehat c_j)\) is the moment
sequence of \(\Phi_1(u)\,du+\tau\delta_0\).  If
\(c_0^{(1)}=\int\Phi_1\), then
\begin{equation}
 \widehat c_0-c_0=y_0,\qquad c_0-c_0^{(1)}=\tau-y_0>0,
\end{equation}
and consequently
\begin{equation}
 \frac{|c_0-c_0^{(1)}|}{|c_0-\widehat c_0|}
 \ge a^2(1-e^{-5\pi})-1>2200.
 \label{eq:green-boundary-gain}
\end{equation}
Here \(\tau\ge\beta_2=ae^{-4\pi}/8\) and
\eqref{eq:green-mode-tail} suffice; \(\pi>3\) and \(e^{15}>1000\)
give the final strict inequality.  This is a gain for the stated
\(c_0\) approximation only, not a determinant-error or runtime gain.

For a finite acquisition cutoff \(U\ge0\), put
\begin{equation}
 a_j^{[U]}=\frac1{8(2j)!}\int_0^Uu^{2j}p_1(u)\,du,
 \qquad a_U=4\pi e^{4U}-1.
\end{equation}
The complete omitted first-mode exterior obeys
\begin{equation}
 0\le x_j-a_j^{[U]}\le B_j(U):=
 \frac{e^{U-\pi e^{4U}}}{8(2j)!}
 \sum_{h=0}^{2j}\binom{2j}{h}
                \frac{U^{2j-h}h!}{a_U^{h+1}}.
 \label{eq:green-spatial-tail}
\end{equation}
To see this, write \(u=U+v\), bound
\(p_1(U+v)\le e^{U-\pi e^{4U}}e^{-a_Uv}\), expand
\((U+v)^{2j}\), and integrate.  This does not pay quadrature uncertainty
on \([0,U]\), which must be enclosed separately.  Nor is a Green
identity applied on the truncated interval.

\subsection{The combined determinant payment}
Fix \(U=2\), write \(a_j=a_j^{[2]}\), and set
\begin{equation}
 T_j=E_j+B_j(2),\qquad
 \mathcal C_j(z)=b\delta_{j0}+a_{j-1}-a_j+z_{j-1}-z_j
 \quad(j\ge0).
\end{equation}
There is an actual vector \(z\) with \(0\le z_j\le T_j\) such that
\(c_j=\mathcal C_j(z)\).  Put \(c_j^*=\mathcal C_j(0)\), and define
\begin{align}
 Q_{r,k}&=D_{r,k}^2-D_{r,k-1}D_{r,k+1}
          =D_{r+1,k}D_{r-1,k},\\
 P_{r,k}(z)&=D_{r,k}(\mathcal C(z))^2
       -D_{r,k-1}(\mathcal C(z))D_{r,k+1}(\mathcal C(z)).
 \label{eq:green-joint-Q}
\end{align}
The product identity is Desnanot--Jacobi, with all boundary zeros retained.
After collecting the joint polynomial, write
\begin{equation}
 P_{r,k}(z)=q^*_{r,k}+\sum_{\alpha\ne0}p_\alpha z^\alpha,
 \qquad \mathcal R_{r,k}=\sum_{\alpha\ne0}|p_\alpha|T^\alpha,
 \qquad L_{r,k}=q^*_{r,k}-\mathcal R_{r,k}.
 \label{eq:green-Taylor-payment}
\end{equation}
Only moments through \(k+r\) and degrees through \(2r\) can occur.
The exact finite expansion immediately proves
\begin{equation}
 |Q_{r,k}-q^*_{r,k}|\le\mathcal R_{r,k},\qquad Q_{r,k}\ge L_{r,k}.
 \label{eq:green-absolute-payment}
\end{equation}
The neighboring products must be subtracted and collected before absolute
values are taken.  The remaining payment still discards correlations
between the actual residual moments; it is sufficient, not necessary.

A useful bound on this same payment avoids listing every translated term.
Let \(H_{n,k}(c^*)\) be the collected determinant polynomial with its
coefficients replaced by their absolute values, evaluated at positive
reference coefficients.  If each nonzero used coefficient satisfies
\(\|\mathcal C_j-c_j^*\|_T\le\epsilon c_j^*\), where
\(\|\sum h_\alpha z^\alpha\|_T=\sum|h_\alpha|T^\alpha\), then
\begin{equation}
 B_{n,k}:=\|D_{n,k}(\mathcal C)-D_{n,k}(c^*)\|_T
 \le[(1+\epsilon)^n-1]H_{n,k}(c^*).
 \label{eq:green-det-norm}
\end{equation}
This follows by expanding each homogeneous degree-\(n\) monomial and
retaining every nonconstant binomial term.  Negative-index entries vanish
identically.  Coefficient convolution proves submultiplicativity of
\(\|\cdot\|_T\).  Thus, with \(d_\pm=D_{r\pm1,k}(c^*)\),
\begin{equation}
 \mathcal R_{r,k}\le
 |d_+|B_{r-1,k}+|d_-|B_{r+1,k}+B_{r+1,k}B_{r-1,k}.
 \label{eq:green-product-payment}
\end{equation}
The last term is the full product remainder.  If \(d_\pm>0\), dividing
by \(q^*_{r,k}=d_+d_-\) gives the relative allowance
\(e_++e_-+e_+e_-\), where \(e_\pm=B_{r\pm1,k}/d_\pm\).
Reference margins and conditioning therefore remain separate inputs.

\subsection{Recorded finite certificates}
\begin{proposition}[Finite quantitative comparison]
For the complete canonical source with the conventions above, the
recorded enclosures and exact arithmetic verification establish
\begin{equation}
 Q_{3,1}>7\cdot10^{-23},\qquad
 Q_{3,4}>10^{-74},\qquad
 Q_{4,1}>3\cdot10^{-31}.
 \label{eq:green-certified-Q}
\end{equation}
For the last case the unchanged compact-reference payment moreover obeys
\begin{equation}
 \frac{\mathcal R_{4,1}}{q^*_{4,1}}<\frac1{15},\qquad
 L_{4,1}>\frac7{3000}(c_1^*)^8
          >\frac7{3000^9}>3\cdot10^{-31}.
 \label{eq:green-next-order-margin}
\end{equation}
\end{proposition}
\begin{proof}[Analytical closure and certificate inputs]
For \((r,k)=(3,1),(3,4)\), the accepted record encloses the eight
compact moments \(a_0,\ldots,a_7\), all mode tails and all spatial tails,
then evaluates \eqref{eq:green-Taylor-payment} outward through degree six.
The serialized dyadic lower allowances exceed the first two rational
floors in \eqref{eq:green-certified-Q}; the saved-data verifier also
recomputes their rational propagation.  These are the finite numerical
inputs, rather than an assumption on the desired determinant sign
(RGD Note~03, the saved E02~r1 record and Review~03;
Appendix~\ref{sec:provenance}).

For the last inequality, exact verification from the same acquisition
supplies \(c_0^*,\ldots,c_5^*>0\), \(c_1^*>1/3000\), the entry-relative
budget \(\epsilon=1/40000\), and the rational boxes
\begin{equation}
 t_2\in[.465,.466],\quad t_3\in[.135,.137],\quad
 t_4\in[.028,.029],\quad t_5\in[.004,.005],\qquad
 t_j=\frac{(c_0^*)^{j-1}c_j^*}{(c_1^*)^j}.
 \label{eq:green-rational-box}
\end{equation}
All terminating decimals here denote exact rationals.  First-row
expansion at \(k=1\) gives
\begin{equation}
 D_{n,1}=\sum_{j=1}^n(-1)^{j-1}c_0^{j-1}c_jD_{n-j,1}.
\end{equation}
Writing \(D_{n,1}(c^*)=(c_1^*)^nf_n\), it follows that
\begin{align}
 f_3&=1-2t_2+t_3,\\
 f_5&=1-4t_2+3t_2^2+(3-2t_2)t_3-2t_4+t_5.
\end{align}
Endpoint monotonicity on \eqref{eq:green-rational-box} gives
\(1/5<f_3<1/4\), \(1/80<f_5<1/40\).  The corresponding collected
absolute polynomials satisfy \(S_3<21/10\), \(S_5<21/5\).
Equation \eqref{eq:green-det-norm}, including all higher binomial terms,
therefore gives
\(B_{3,1}/(c_1^*)^3<21/100000\) and
\(B_{5,1}/(c_1^*)^5<63/100000\).  Substitution in
\eqref{eq:green-product-payment} yields
\begin{equation}
 \frac{\mathcal R_{4,1}}{(c_1^*)^8}
 <\frac1{40}\frac{21}{100000}
  +\frac14\frac{63}{100000}
  +\frac{21\cdot63}{10^{10}}
 =\frac{1628823}{10^{10}}<\frac1{6000},
\end{equation}
whereas \(q^*_{4,1}>(c_1^*)^8/400\).  Subtraction proves
\eqref{eq:green-next-order-margin}.  The second exact record verifies
these input-box memberships, the polynomial identities and every rational
inequality; it acquires no additional integral
(RGD Notes~05--06, the saved E02~r2 record and Review~04;
Appendix~\ref{sec:provenance}).
\end{proof}

The first acquisition used python-flint~0.9.0 and FLINT~3.6.0 at
384-bit precision, retaining forty segment enclosures for eight moments.
The exact saved-data verifiers trust those identified integral and
constant enclosures; they do not independently reproduce the integration
or formally verify FLINT.  The analytical comparison received
cross-model internal LLM assessment.  Assessment of the first acquisition
was by its originating assistant/model with fresh same-model contributions;
the next-order argument received cross-model checking, while its new
implementation retained author checks.  None is independent specialist
review (RGD Reviews~01, 03 and~04; Appendix~\ref{sec:provenance}).

The certified quantities are the products \(Q_{r,k}\) defined in
\eqref{eq:green-joint-Q}.  No numerical normalization by a lower-order
determinant or relative induction parameter was acquired.  We assert no
new positivity range, and these finite bounds do not supply an all-order
or all-index conclusion.

\subsection{Why the successful finite payment cannot be iterated unchanged}
For the specified compact reference and independent moment bounds,
there is an exact obstruction stronger than poor conditioning of an
auxiliary estimate.  Let
\(K_0=e^{-4\pi}/[8(1-e^{-5\pi})]\), so \(E_j=K_0/a^{2j+1}\).
In the already collected joint polynomial,
\begin{equation}
 [z_{k-1}^{2r}]P_{r,k}=1,\qquad
 \mathcal R_{r,k}\ge E_{k-1}^{2r}.
\end{equation}
Indeed, its homogeneous dependence on \(z_{k-1}=v\) gives coefficient
changes \(\mathcal C_{k-1}=-v\), \(\mathcal C_k=v\); each relevant
index-\(k\) matrix is lower bidiagonal, with determinant \(v^n\).
For \(j\ge2\), the compact support gives
\(0<c_j^*\le a_0 2^{2j-2}/(2j-2)!\).  If \(k\ge r+2\),
log-convexity of the factorial along each permutation consequently gives
\begin{equation}
 |q^*_{r,k}|\le(r+1)!(r-1)!
       \left(\frac{a_0 2^{2k-2}}{(2k-2)!}\right)^{2r}.
\end{equation}
The ratio of \(E_{k-1}\) to the parenthesized quantity is
\begin{equation}
 A_k=\frac{K_0}{aa_0}\frac{(2k-2)!}{(2a)^{2k-2}},\qquad
 \frac{A_{k+1}}{A_k}=\frac{(2k)(2k-1)}{(2a)^2}\longrightarrow\infty.
\end{equation}
Thus, for every fixed \(r\ge2\),
\begin{equation}
 L_{r,k}<0\qquad\text{for all sufficiently large }k.
 \label{eq:green-exact-payment-stop}
\end{equation}
This excludes this exact compact-reference Taylor payment, not the actual
canonical determinant or a correlated source residual.  In particular,
removing the spatial cutoff changes the error problem: Note~07 proves
\(0<(c_j-c_j^{[1]})/c_j^{[1]}\le128e^{-j/(16\log j)}\) for
\(j\ge e^{16}\).  Its reference signs and conditioning still require
separate control (RGD Note~07; Appendix~\ref{sec:provenance}).  The finite certificates above
therefore furnish no all-index induction contract.
